\documentclass[11pt]{article}
\usepackage[T1]{fontenc}
\usepackage[utf8]{inputenc}
\usepackage{amsmath, amssymb}
\usepackage{booktabs}
\usepackage{graphicx}
\usepackage{geometry}
\geometry{margin=2.4cm}
\usepackage{setspace}
\usepackage[bottom]{footmisc}
\usepackage[hidelinks]{hyperref}
\usepackage{threeparttable}
\usepackage{float}
\usepackage[section]{placeins}
\usepackage{caption}
\newcommand{\sym}[1]{$^{#1}$}
\graphicspath{{/Users/edoardovitella/Documents/work/projects/Paper_PTA/New/Paper/paper_v3/figures/}}

% --- placeholder refs so the draft compiles standalone -----------------------
\newcommand{\phlabel}[1]{\phantom{\label{#1}}}
% lightweight citation stubs (no bibliography in this standalone draft)
\newcommand{\citep}[1]{(\textsc{#1})}
\newcommand{\citet}[1]{\textsc{#1}}

\title{\textbf{Results --- proposed reorganisation (v4 draft)}\\[6pt]
\large Environmental Provisions in Chinese PTAs and Export Composition}
\author{Edoardo Vitella}
\date{\today}

\begin{document}
\maketitle

\noindent\rule{\textwidth}{0.4pt}
\footnotesize
\textbf{What this document is.} A standalone draft of a reorganised Section~5. Six subsections in
place of twelve, ordered: full panel $\rightarrow$ collapsed panel and inference $\rightarrow$
control groups $\rightarrow$ dynamics $\rightarrow$ sub-indices $\rightarrow$ extensive margin.
All figures come from the corrected table document; the numbers that changed relative to
\texttt{paper\_v3} are flagged in footnotes. Cross-references to appendix tables are
placeholders. Nothing in \texttt{paper\_v3.tex} has been modified.
\normalsize
\noindent\rule{\textwidth}{0.4pt}

\vspace{12pt}
\noindent\emph{Specification recap.} The estimating equation is
\begin{equation}\label{eq:main}
\ln x_{fpdt} = \beta_1\, EP_{dt}\!\times\!g_p + \beta_2\, EP_{dt}\!\times\!b_p
+ \gamma_1\, TD_{dt}\!\times\!g_p + \gamma_2\, TD_{dt}\!\times\!b_p
+ \theta_{fpd} + \theta_{fdt} + \theta_{pt} + \varepsilon_{fpdt},
\end{equation}
with $g_p$ and $b_p$ the green and pollution-intensive product indicators, $EP_{dt}$ environmental
depth and $TD_{dt}$ non-environmental agreement depth. Standard errors are clustered by destination
throughout; because the design has 23 treated destinations, the wild cluster bootstrap
($B=9{,}999$) rather than the asymptotic $p$-value is the reference inference.

% ============================================================================
\section{Results}\label{sec:results}
% ============================================================================

\subsection{The full panel}\label{sec:fullpanel}

Table~\ref{tab:fullpanel} reports the reference estimates. Equation~\eqref{eq:main} is estimated on
the full firm--product--destination--year panel: 21.5 million observations, 225 destination
clusters, firm$\times$product$\times$destination, firm$\times$destination$\times$year and
product$\times$year fixed effects. Columns~(1)--(3) use the World Bank EP depth index, columns
(4)--(6) the TREND provision count; within each block the outcome is log export value, log
quantity and log unit value. Panel~A controls for non-environmental World Bank depth, Panel~B for
the independently coded DESTA index \citep{dur2014}.

The table is deliberately wide, because the central claim of this section is a claim about
\emph{invariance}: the green coefficient does not move across any of these dimensions, and the
dirty coefficient moves across exactly one of them.

\paragraph{The green margin.} Zero everywhere. On export value the WB estimate is $-0.0023$
(s.e.\ 0.0039, bootstrap $p = 0.69$) under the World Bank depth control and $-0.0019$
(0.0037, $p = 0.67$) under DESTA; the TREND estimate is $-0.0001$ in both, with bootstrap
$p$-values of 0.93 and 0.92. The joint $F$-test on all four composition interactions gives
$p = 0.31$ (WB) and $p = 0.71$ (TREND) in Panel~A, and $p = 0.18$ and $p = 0.26$ in Panel~B.
Quantity and unit value behave identically: across the eight green cells in columns~(2), (3),
(5) and~(6), no bootstrap $p$-value falls below 0.74.

That last point does work that the value regression alone cannot do. A null on export value is
consistent with two offsetting responses --- more units at lower prices, or fewer units at higher
prices --- cancelling in the aggregate. Columns~(2) and~(3) rule this out directly: neither the
volume channel nor the price channel moves. The green null is a null on both margins separately,
not a coincidence of the value aggregate.\footnote{This replaces the summary table in
\texttt{paper\_v3} \S5.10, which reported ranges of $p$-values and noted that exact coefficients
were ``available upon request''. They are now in columns~(2), (3), (5) and~(6).}

\paragraph{The dirty margin.} Negative throughout, and this is the one place in the paper where a
modelling choice changes the verdict. Under the World Bank depth control the WB$\times$dirty
coefficient is $-0.0044$ with asymptotic $p = 0.052$, but the bootstrap raises it to $p = 0.19$;
the TREND counterpart is $-0.0009$ ($0.15$ asymptotic, $0.18$ bootstrap). Under the DESTA control
the same coefficients are $-0.0056$ ($0.039$ asymptotic, \textbf{$0.049$ bootstrap}) and $-0.0014$
($0.048$ asymptotic, $0.069$ bootstrap). Same data, same fixed effects, same estimation sample; the
only difference is which measure of non-environmental agreement depth is partialled out.\footnote{%
In \texttt{paper\_v3} this contrast is invisible, because the tables reported only the World Bank
depth control in the main text and the sentence ``the TREND index never shows the effect'' was
written from the collapsed-panel figure. On the full panel with the DESTA control the TREND index
does show it. The corrected tables make the contrast the substantive question it is.}

The mechanical explanation is collinearity. Net of the fixed effects, environmental depth
correlates $0.959$ with World Bank non-environmental depth and $0.891$ with DESTA depth. The more
the control overlaps the regressor of interest, the less independent variation is left and the
wider the standard error. Read that way, Panel~B is the better-identified specification and Panel~A
is over-controlled --- which would make the dirty effect real and the World Bank specification the
one hiding it.

We do not adopt that reading, and Section~\ref{sec:collapsed} explains why: the permutation test,
which reassigns entire EP profiles across treated destinations and is therefore indifferent to
which depth control is in the regression, assigns this coefficient a $p$-value between 0.15 and
0.38 in every one of the four sample--control variants. But the ambiguity is real and we flag it
here rather than resolving it by choosing the control that gives the answer we prefer.

\paragraph{Quantity, not price.} Where the dirty coefficient does move, it moves on quantity.
Under the DESTA control it is $-0.0070$ on log quantity (asymptotic $p = 0.020$, bootstrap
$p = 0.035$) --- larger in magnitude and more precisely estimated than on value --- while the unit
value coefficient is $+0.0004$ with $p = 0.68$. Taking Tables~\ref{tab:fullpanel} and
\ref{tab:collapsed} together, sixteen cells estimate an EP interaction on log unit value; exactly
one is nominally significant (TREND$\times$dirty in the collapsed panel, asymptotic $p = 0.076$)
and it does not survive the bootstrap ($p = 0.17$). Whatever this coefficient reflects, it is a
shipment-volume margin with no price response. That is consistent with a compositional shift in
what firms ship and inconsistent with a quality-upgrading or cost-passthrough interpretation.

\paragraph{How large a null.} A null result means one of two things: the effect is absent, or the
design cannot see it. The bootstrap confidence intervals say which. On the full panel the 95\%
bootstrap interval for EP$\times$green is $[-0.0353,\, +0.0355]$ per World Bank provision --- about
$\pm 3.5\%$, or, since one standard deviation of WB EP depth in the estimation sample is roughly
2.4 provisions, at most a $\pm 8$--$9\%$ change in green relative to neutral exports. The
corresponding asymptotic interval, $[-0.0100,\, +0.0055]$, is roughly six times narrower; reporting
that width would claim a precision this paper's own inference section declares unwarranted.

Table~\ref{tab:brandi} places the bound against the closest benchmark. \citet{brandi2020} find that
one liberal environmental provision raises the green share of developing-country exports by around
0.4 percentage points, which is roughly $+0.16$ log points in the metric of
equation~\eqref{eq:main}. The upper end of the interval here is about one quarter of that figure.
The design can therefore exclude effects at or above the Brandi magnitude; it cannot exclude
smaller positive effects. Appendix Table~\ref{tab:mde} reports the implied minimum detectable
effects in per-provision and per-standard-deviation terms.

\input{tabs/tab_fullpanel}

\input{tabs/tab_brandi}

\FloatBarrier
\subsection{The collapsed panel, and what survives robust inference}\label{sec:collapsed}

Table~\ref{tab:collapsed} re-estimates equation~\eqref{eq:main} on HS6--destination--year cells with
analytic weights. This is not a robustness check on Table~\ref{tab:fullpanel}. It estimates a
different object, and the difference between the two is itself a result.

The firm$\times$destination$\times$year fixed effects in the full panel absorb everything that
varies at the firm--destination--year level, so identification comes from reallocation
\emph{within} the firm across products. The collapsed panel compares products across all firms
shipping to a destination in a year, so any shift in \emph{which} firms export dirty goods enters
the coefficient as well. The two estimands differ by exactly the between-firm channel.

On the green margin they agree: $-0.0046$ collapsed against $-0.0023$ full panel, both
indistinguishable from zero on every inference method. On the dirty margin they differ by a factor
of roughly 2.7 ($-0.0119$ against $-0.0044$), which decomposes the coefficient: approximately
three-fifths of the collapsed dirty estimate, $(0.0119 - 0.0044)/0.0119 = 0.63$, is between-firm
composition rather than within-firm reallocation. Whatever moves dirty exports operates mainly
through which firms ship, not through what a given firm ships. The full panel therefore remains the
reference specification, and the collapsed panel is where the between-firm channel becomes visible.

\input{tabs/tab_collapsed}

\paragraph{Why inference lives here.} The collapsed panel is also where the design can be stress
tested. The permutation test reassigns the EP profile of each agreement --- its depth level and its
entry years --- randomly among the 23 treated destinations and re-estimates the model a thousand
times; the treated-only design this requires is not feasible at firm level. The leave-one-out
exercise, likewise, is a collapsed-panel exercise. Table~\ref{tab:collapsed} therefore carries three
rows of inference where Table~\ref{tab:fullpanel} carries two.

Read down those rows for the WB$\times$dirty coefficient. The asymptotic $p$-value is below 0.001.
The wild cluster bootstrap raises it to 0.072. The permutation test raises it to 0.278: in 27.8\% of
random reassignments of environmental content across the same 23 destinations, the model produces a
coefficient at least as large in absolute value. A coefficient that would have arisen a quarter of
the time from randomly relabelling which agreements are environmental is not evidence about
environmental content.

\paragraph{The depth-control ambiguity, resolved as far as it can be.} Table~\ref{tab:matrice}
assembles the dirty coefficient across all four sample--control variants with a single inference
method, so the eight cells are directly comparable. Two features stand out. The point estimate is
uniformly negative and, within each panel, stable. The bootstrap $p$-value is not: on the collapsed
panel it is 0.072 with the World Bank control and 0.047 with DESTA; on the full panel it is 0.185
and 0.049. Significance tracks the control variable, not the data.

The permutation test is the tiebreaker, because it is the one procedure whose null distribution
does not depend on which depth control sits in the regression. It gives $p = 0.278$, $0.146$,
$0.146$ and $0.384$ across the four variants --- no cell below 0.10. The reading we adopt is
therefore not that the dirty coefficient is a false positive, and not that it is an established
effect, but that it is a coefficient whose apparent precision is an artefact of how much
independent variation a particular depth control leaves behind. On the evidence that does not
depend on that choice, it is indistinguishable from a placebo.

\input{tabs/tab_matrice}

\paragraph{Leave-one-out: the mechanism is precision, not leverage.} Table~\ref{tab:loo_new}
re-estimates the dirty coefficient dropping one treated destination at a time. The usual concern
--- that a result is driven by a single country --- is not what the exercise finds. The point
estimate sits between $-0.0097$ and $-0.0133$ in all 23 leave-one-out samples; dropping large
treated destinations such as India or Pakistan moves it by one to four per cent and leaves the
standard error essentially unchanged.

What is unstable is precision. Dropping Australia moves the coefficient by 13\%, to $-0.0103$, but
nearly triples its standard error, from 0.0030 to 0.0087, carrying the $p$-value to 0.24. Dropping
South Korea doubles the standard error, giving $-0.0097$ with $p = 0.09$. Australia and Korea are
not outliers exerting leverage on the estimate; they supply the identifying variation, and removing
them leaves too little information for a precise estimate rather than revealing a different effect.
The coefficient rests on a thin slice of variation concentrated in two destinations --- and which
two depends on the depth control, since under DESTA the pivotal destination shifts from Australia
to Korea. That a modelling choice unrelated to environmental content determines which destination is
pivotal is the same fragility diagnosis reached in Table~\ref{tab:matrice} by a different route.

One caveat we do not resolve. Australia and Korea are also among the destinations with the deepest
EP coverage in the sample (WB depth 12 and 17). A dose--response reading --- the dirty effect is
real and concentrated where the dose is largest --- is not ruled out by the leave-one-out geometry.
The row \textit{High EP depth only} in Table~\ref{tab:loo_new}, which drops Peru, Switzerland and
Korea jointly, is consistent with either interpretation.

\input{tabs/tab_loo_new}

\FloatBarrier
\subsection{Control groups and estimation designs}\label{sec:subsamples}

The estimates so far compare green and dirty products against neutral products, in agreement
destinations against non-agreement destinations. Both comparisons can be challenged. Neutral
products may be poor counterfactuals for green ones if they sit in different parts of the product
space; non-agreement destinations may be poor counterfactuals for agreement destinations if China
signs agreements with countries on different growth paths. Table~\ref{tab:subsamples} replaces each
comparison in turn.

\emph{C-prod-HS4} restricts attention to HS6 products with a comparable HS4-level match, so a green
product is compared with neutral products in the same four-digit heading rather than with the whole
neutral universe. \emph{CEM} matches destinations on log GDP per capita, growth and MFN tariff in
2000, reducing the sample to 57 destination clusters that are comparable on pre-treatment
observables. \emph{Deep versus shallow} discards untreated destinations altogether: identification
comes from variation in EP depth among the 23 PTA partners, so a destination with a shallow
environmental chapter is the counterfactual for one with a deep chapter.

Every cell in the table uses the wild cluster bootstrap, so all forty are comparable, and no stars
are reported: with this many designs, star-counting invites exactly the selective reading the table
is meant to prevent.

The green coefficient stays between $-0.0046$ and $+0.0005$ in all twenty cells and is never close
to conventional significance --- the smallest bootstrap $p$-value anywhere on the green margin is
0.51. This is the result. An artifact of the comparison group would move when the comparison group
changes; this does not move when the product comparison is narrowed to within-HS4, when the
destination comparison is matched on pre-treatment observables, or when untreated destinations are
removed from the design entirely.

The dirty coefficient behaves differently, and in a way that reinforces
Section~\ref{sec:collapsed}. It is negative in eight of the ten firm-level cells but never
significant once the sample changes: CEM gives $-0.0041$ ($p = 0.25$) and $-0.0043$ ($p = 0.12$),
deep-versus-shallow gives $-0.0034$ ($p = 0.50$) and $-0.0042$ ($p = 0.17$). C-prod-HS4 reverses
the sign entirely, to $+0.0066$ and $+0.0033$. A coefficient that survives the removal of the
untreated comparison group only in the specification that keeps it is not identifying a robust
composition effect.

\input{tab_subsamples}

Further sample variations --- adding product--destination--year controls for MFN tariffs, market
concentration and antidumping exposure; excluding the eleven ASEAN destinations, which form the
largest single cohort; re-including Hong Kong and Macao; and restricting to products with common
support across agreement and non-agreement destinations --- are reported in Appendix
Table~\ref{tab:robust-full}. None moves the green coefficient outside $[-0.0025, -0.0002]$.
Appendix Table~\ref{tab:apec} replaces the OECD list of environmental goods with the narrower APEC
list agreed in 2012; the green coefficient changes sign but not significance, confirming that the
null does not depend on which products are labelled green.

\FloatBarrier
\subsection{Dynamics and pre-trends}\label{sec:dynamics}

Figure~\ref{fig:es} plots the event study of the green and dirty gap relative to neutral products
around entry into force, using never-treated destinations as controls. Differential pre-trends are
flat on both margins.

Before reading that as reassurance, a clarification. Parallel trends is a claim about the
unobserved counterfactual path of treated destinations \emph{after} treatment. A flat pre-trend
rules out the most easily detectable violation --- a composition gap already in motion before the
agreement --- but does not establish that the assumption holds thereafter. The exercise is a
falsification, not a proof.

\begin{figure}[htbp]
\centering
\caption{Event study: green and dirty vs.\ neutral products around PTA entry into force}
\label{fig:es}
\vspace{6pt}
\includegraphics[width=\textwidth]{eventstudy_collapsed_v3.png}
\begin{minipage}{\textwidth}\footnotesize
\emph{Notes:} Collapsed panel; FE $pd + dt + pt$; destination-clustered 95\% confidence bands;
$t = -1$ is the reference period; endpoint bins accumulate. Coefficients are tabulated in Appendix
Table~\ref{tab:eventstudy_twfe}. Source: Chinese Customs Office; World Bank DTA Database.
\end{minipage}
\end{figure}

A mild negative drift for green products in the late post-period (bins $\geq +5$) is identified
only by the early cohorts --- ASEAN 2005, Chile 2006, Pakistan 2007, New Zealand 2008. It disappears
under the \citet{sun2021} interaction-weighted estimator applied to the destination-level
composition gap: the aggregated ATT is $-0.042$ ($p = 0.27$) for green and $+0.073$ ($p = 0.28$)
for dirty, consistent with early-cohort heterogeneity rather than any effect of treatment. The
Sun--Abraham estimator answers a narrower question than equation~\eqref{eq:main} --- it binarises
treatment, discards depth, carries no depth control, and operates on a gap collapsed across
products --- so it is a diagnostic for staggered timing, not a replication. The full
interaction-weighted event study is in Appendix Table~\ref{tab:sunab}; with standard errors that
account for estimation uncertainty in the cohort-share weights, no coefficient in the $[-10, +8]$
window is individually distinguishable from zero on either margin.

\paragraph{Destination-specific trends.} The fixed effects in equation~\eqref{eq:main} do not
absorb a shock specific to green products in a particular destination that grows over time --- a
country signing environmental clauses precisely as its demand for green goods is rising. Two
variants bracket the answer, both on the collapsed panel (Table~\ref{tab:trends}).

Allowing each destination its own linear trend interacted with product type over the full sample
period leaves the World Bank results essentially unchanged (EP$\times$green $-0.0051$, bootstrap
$p = 0.49$; EP$\times$dirty $-0.0082$, bootstrap $p = 0.27$). The TREND green interaction, however,
turns significantly negative and is the single coefficient in the entire robustness battery that
survives the wild cluster bootstrap ($p = 0.015$).

This should not be taken at face value, for two reasons. First, destination-specific trends
estimated over the full sample period absorb both pre- and post-treatment variation and can reverse
the sign of treatment coefficients when cohort dynamics are present \citep{wolfers2006}; the late
negative drift these trends absorb is precisely the early-cohort artefact the Sun--Abraham
estimator identifies as noise. Second, and more directly, Panel~B$'$ shows the green gap was already
rising before entry into force (pre-agreement slope $+0.0073$, bootstrap $p = 0.19$), so a
pre-existing trend is a plausible account of the same coefficient.

The second variant estimates destination-specific slopes on pre-treatment years only and projects
them forward before re-estimating on the detrended outcome. Every coefficient returns to an
imprecise zero: TREND$\times$green $+0.0073$ ($p = 0.19$), WB$\times$green $+0.0173$ ($p = 0.68$),
dirty interactions $p = 0.29$--$0.61$. The composition null stands against destination-specific
linear drifts in green demand, which is the most natural form of the selection-on-trends
confounder.

\input{tabs/tab_trends}

\FloatBarrier
\subsection{Provision bundling and the limits of heterogeneity analysis}\label{sec:bundling}

The indices used so far sum all environmental provisions in an agreement regardless of content.
Most are cooperation clauses, institutional commitments or general environmental language, none of
which has an identifiable channel through which it could shift which products a country exports.
Only a small minority carry a direct trade mechanism: tariff reductions on green goods,
non-regression clauses preventing countries from lowering standards to attract trade, or binding
market-access commitments. A null on the aggregate index is therefore ambiguous. It could mean that
environmental provisions do not affect export composition, or that the signal from a few
trade-relevant clauses is diluted by a large majority that should have no effect in the first place.

Table~\ref{tab:mechanism} quantifies the dilution. The World Bank index checks 150 provision slots
across the 25 treated destinations; 8 of them (5.3\%) belong to sub-indices with a trade mechanism.
The TREND database, coded at finer granularity, gives 4 of 437 (0.9\%). At agreement level, a
binding environmental obligation appears in exactly one of the 14--15 Chinese PTAs in the database:
China--Korea, 2015.

\input{tabs/tab_mechanism}

The structural consequence is stark. The two World Bank sub-indices with a direct trade mechanism
are non-zero in only three destination--years, always in the same proportion, and are therefore
perfectly collinear ($\rho = 1.000$) and not separately identifiable. \citet{brandi2020} report a
maximum VIF of 4.6 across provision-type variables on a 680-agreement panel; with 13 distinct EP
profiles, this design cannot do the same.

A reading rule applies to everything that follows. Each sub-index replaces the aggregate index one
at a time and is estimated on its own, so the coefficients are not partial effects holding other
clause types constant --- they absorb whatever correlated variation moves with each one.
Table~\ref{tab:subindex} is a set of separate comparisons, not an additive decomposition.

Two things can be said. First, the trade-mechanism bundle reproduces the aggregate pattern: green
never responds (WB Green Liberalization$\times$green $p = 0.90$; TREND Green Market
Access$\times$green $p = 0.37$), while the dirty interaction is marginally negative on asymptotic
$p$-values --- the same asymptotics-only signal documented in Section~\ref{sec:collapsed}.

Second, and more usefully, the full panel and the collapsed panel disagree on the sub-indices in
exactly the way they disagree on the aggregate. On the collapsed panel, WB Green
Liberalization$\times$dirty is $-0.088$ and TREND Green Market Access$\times$dirty is $-0.049$; on
the full panel the same coefficients are $-0.002$ and $-0.009$ and nowhere near significance
(Appendix Table~\ref{tab:subindex_fullpanel}). The sub-index results, like the aggregate dirty
coefficient, are a between-firm phenomenon that does not survive firm$\times$destination$\times$year
fixed effects.

The placebo sub-indices deserve care. TREND soft provisions show nothing on either margin
($\times$green $p = 0.27$; $\times$dirty $p = 0.86$), consistent with the logic of the null.
Regulatory-space clauses do not: they carry $+0.024$ on green and $+0.023$ on dirty, and both
survive the wild cluster bootstrap ($p = 0.045$ and $p = 0.021$). This is the only place in the
paper where a sub-index without a trade mechanism registers a signal that robust inference sustains.

Two features bound what can be read into it. The two margins move in tandem --- the green--dirty
differential is $+0.0017$ with $p = 0.80$ --- so the coefficient describes green and dirty exports
rising equally against the neutral baseline, which is not a shift in composition toward cleaner
goods. And the regulatory-space sub-index is not a clean placebo: it accounts for 71.5\% of the
TREND count across treated destination--years and correlates 0.90 with total depth, and in this
specification the depth control flips to significantly negative. Two regressors correlated at 0.90
splitting into a large positive and a large negative is the bundling problem applied to the
placebo, not a cleanly identified channel.

\input{tabs/tab_subindex}

This reconciles the null with \citet{brandi2020}. Their finding is not that environmental provisions
generically matter; it is driven by the subset of agreements with explicitly liberalizing or
restrictive clauses, identified separately because 680 agreements supply enough cross-agreement
variation in provision type. The Chinese PTAs of this period are, on the same coding, overwhelmingly
cooperation-only. A null on the aggregate index is what their own heterogeneity would predict for a
sample of this composition. A second difference compounds it: \citet{brandi2020} work with
exporter--importer--year flows, where a composition effect can arise from the entry or exit of firms
or products, whereas equation~\eqref{eq:main} identifies from reallocation within the same firm ---
a strictly more demanding test of the same hypothesis.

\FloatBarrier
\subsection{The extensive margin}\label{sec:ppml}

Log-linear estimates condition on positive flows. If environmental provisions create new green trade
where none previously existed, the estimates above miss it \citep{santossilva2006}. On an
HS6--destination--year grid completed with zeros (8.2 million cells, of which 7.9 million enter
estimation after PPML drops separated and singleton observations), the picture is unchanged
(Table~\ref{tab:ppml}): EP$\times$green is $+0.0015$ ($p = 0.74$, WB) and $+0.0001$ ($p = 0.95$,
TREND), and EP$\times$dirty is $-0.030$ ($p = 0.057$, WB) and $+0.003$ ($p = 0.55$, TREND). No green
trade creation is detected at the extensive margin, and the dirty coefficient is at most marginal,
consistent with Section~\ref{sec:collapsed}.

One limitation should be noted. This grid detects a green product newly exported to a destination,
but not a new firm beginning to export a product that other Chinese firms already ship there, which
is a common form of firm-level entry. The extensive margin is probed, not exhausted.

\input{tabs/tab_ppml}

\FloatBarrier
\vspace{18pt}
\noindent\rule{\textwidth}{0.4pt}
\subsection*{Note on what moved to the appendix}
\footnotesize
Saturation ladder (with a baseline column to be added, so the sequence terminates at
Table~\ref{tab:fullpanel}); Sun--Abraham event study; tabulated TWFE event study; Hong Kong and
Macao variants of both main tables; full wild-cluster-bootstrap and permutation tables; full-panel
robustness battery; APEC versus OECD product definition; CO$_2$ intensity as a continuous dirty
measure; sub-index estimates on the full panel; minimum detectable effects; environmental share of
agreement content \citep{abman2024}; the full subsample tables underlying
Table~\ref{tab:subsamples}.

\bigskip
\noindent\textbf{Open drafting decisions.}
(i)~Table~\ref{tab:trends} is included whole; Panels~A/A$'$ could move to the appendix leaving only
the pre-trend test in the main text. (ii)~The saturation-ladder sign reversal on the dirty margin
($+0.0064$ under $fpd + pt$, $-0.0044$ in the reference specification) is a strong fragility
argument and currently sits only in the appendix; it could be a paragraph in
Section~\ref{sec:collapsed}. (iii)~Table~\ref{tab:mde} is estimated on the collapsed panel while the
bounding paragraph in Section~\ref{sec:fullpanel} uses full-panel intervals; these should be
harmonised.
\normalsize

\newpage
\appendix
\section*{Appendix placeholders (labels only, tables live in the table document)}
\footnotesize
\begin{itemize}\itemsep2pt
\item Saturation ladder \phantomsection\label{tab:ladder_wb}\label{tab:ladder_trend}
\item Sun--Abraham event study \phantomsection\label{tab:sunab}
\item Event study (TWFE), tabulated \phantomsection\label{tab:eventstudy_twfe}
\item Full panel incl.\ HK/Macao \phantomsection\label{tab:fullpanel_inclHKMO}
\item Collapsed panel incl.\ HK/Macao \phantomsection\label{tab:collapsed_inclHKMO}
\item Wild cluster bootstrap, all variants \phantomsection\label{tab:wcb}
\item Permutation test, all variants \phantomsection\label{tab:perm}
\item Full-panel robustness battery \phantomsection\label{tab:robust-full}
\item APEC vs.\ OECD green definition \phantomsection\label{tab:apec}
\item CO$_2$ intensity (Shapiro) \phantomsection\label{tab:co2}
\item Sub-indices, full panel \phantomsection\label{tab:subindex_fullpanel}
\item Minimum detectable effects \phantomsection\label{tab:mde}
\item C-prod-HS4, full table \phantomsection\label{tab:cprodhs4}
\item CEM destinations, full table \phantomsection\label{tab:cem}
\item Deep vs.\ shallow, full table \phantomsection\label{tab:deepshallow}
\end{itemize}
\normalsize

\end{document}
